\chapter{Meningeal Vein Thrombosis}
\index{Meningeal Vein Thrombosis}

\section*{Definition and Overview}
Meningeal Vein Thrombosis (also referred to under the broader clinical spectrum of cerebral venous thrombosis or cortical vein thrombosis) is a rare but serious neurological condition characterised by the formation of a blood clot (thrombus) within the meningeal or cerebral veins. The occlusion of these venous channels disrupts the normal drainage of blood from the brain tissue, leading to venous congestion, localised cerebral oedema, and elevated intracranial pressure. If untreated, it can cause venous infarction, intracranial haemorrhage, and permanent neurological deficits. Meningeal vein thrombosis is clinically challenging to diagnose because of its highly variable presentation, often mimicking other neurological disorders. Prompt diagnosis using advanced neuroimaging and early initiation of anticoagulant therapy are crucial to prevent catastrophic progression and ensure favourable clinical recovery.

\section*{Epidemiology}
Cerebral venous and meningeal vein thrombosis is an uncommon condition, accounting for approximately 0.5\% to 1\% of all stroke cases globally. It predominantly affects young adults, with a peak incidence in the third and fourth decades of life. There is a marked female-to-male predisposition, with females accounting for approximately 75\% of cases; this gender asymmetry is strongly linked to gender-specific risk factors such as oral contraceptive use, pregnancy, the postpartum period, and hormone replacement therapy. In many countries and other developed nations, the annual incidence is estimated at 1.3 to 1.5 cases per 100,000 population. Paediatric cases also occur, often associated with local parameningeal infections (e.g. otitis media or mastoiditis) or underlying congenital prothrombotic states.

\section*{Aetiology and Pathophysiology}
The thrombogenic process in the meningeal veins is driven by the classic Virchow's triad: endothelial injury, stasis of blood flow, and systemic hypercoagulability.
\begin{itemize}
    \item \textbf{Transient hormonal factors:} Elevated oestrogen levels from combined oral contraceptives, pregnancy, or the postpartum state, which increase hepatic synthesis of clotting factors.
    \item \textbf{Inherited thrombophilias:} Congenital coagulation disorders such as Factor V Leiden mutation, prothrombin gene mutation (G20210A), or deficiencies in protein C, protein S, or antithrombin III.
    \item \textbf{Acquired hypercoagulable states:} Autoimmune conditions, most notably antiphospholipid syndrome (APS), systemic lupus erythematosus, or active malignancies.
    \item \textbf{Parameningeal infections:} Spread of localised bacterial infections from the middle ear (otitis media), mastoid air cells (mastoiditis), or paranasal sinuses, inducing septic thrombus formation.
    \item \textbf{Local mechanical factors:} Closed head trauma, neurosurgical procedures, or lumbar puncture that disrupt venous endothelial integrity or alter cerebrospinal fluid dynamics.
\end{itemize}

\section*{Clinical Features}
The clinical presentation is highly variable, ranging from isolated headache to severe focal neurological deficits and altered consciousness.
\begin{itemize}
    \item \textbf{Intractable headache:} The most common and often earliest symptom (present in up to 90\% of cases), typically progressive, severe, and worsening with coughing, straining, or lying flat (indicative of raised intracranial pressure).
    \item \textbf{Focal neurological deficits:} Hemiparesis, sensory loss, or aphasia, occurring due to localised venous infarction or lobar haemorrhage.
    \item \textbf{Seizures:} Focal or generalised seizures, which are significantly more common in venous thrombosis (up to 40\% of cases) than in arterial ischaemic strokes.
    \item \textbf{Visual disturbances:} Blurry vision, double vision (diplopia due to sixth cranial nerve palsy), or papilloedema (swelling of the optic disc) on ophthalmoscopic exam.
    \item \textbf{Encephalopathy:} Altered mental status, confusion, stupor, or coma in patients with extensive bilateral venous occlusion.
\end{itemize}

\section*{Differential Diagnosis}
Meningeal vein thrombosis must be distinguished from other acute neurological and infectious conditions.
\begin{itemize}
    \item \textbf{Arterial ischaemic stroke:} Acute focal neurological deficits with a rapid onset (seconds to minutes) and vascular distribution, typically lacking the progressive headache and seizures characteristic of venous thrombosis.
    \item \textbf{Bacterial or viral meningitis:} Presenting with fever, headache, photophobia, and neck stiffness; cerebrospinal fluid analysis is diagnostic.
    \item \textbf{Intracranial abscess or tumour:} Mass lesions causing raised intracranial pressure and focal deficits, distinguished by characteristic ring enhancement or tissue architecture on contrast MRI.
    \item \textbf{Subarachnoid haemorrhage:} Classically presenting as a sudden-onset, excruciating "thunderclap" headache, distinguished on non-contrast CT.
    \item \textbf{Pseudotumour cerebri (idiopathic intracranial hypertension):} Presenting with headache, papilloedema, and elevated opening pressure on lumbar puncture, but without venous thrombus on MR venography.
\end{itemize}

\section*{When to Seek Medical Care}
Patients presenting with a progressive, unusually severe headache, new-onset seizures, or focal neurological signs must seek immediate emergency medical evaluation.

\textbf{Contact your local emergency services for:}
\begin{itemize}
    \item \textbf{Sudden severe headache:} A "thunderclap" headache or an intense headache accompanied by neck stiffness or vomiting.
    \item \textbf{Acute neurological deficits:} Sudden weakness or numbness on one side of the face or body, difficulty speaking, or loss of vision.
    \item \textbf{Seizure activity:} New-onset focal or generalised convulsions.
    \item \textbf{Altered consciousness:} Loss of consciousness, severe confusion, or difficulty waking up.
\end{itemize}

\section*{Diagnosis and Investigations}
Diagnosis relies on high-resolution neuroimaging to visualise the venous thrombus and evaluate brain parenchyma.
\begin{itemize}
    \item \textbf{MRI and Magnetic Resonance Venography (MRV):} The most sensitive imaging combination, demonstrating the absence of flow in the affected vein or sinus and showing direct clot visualisation (the "empty delta sign" on contrast-enhanced scans).
    \item \textbf{CT Venography (CTV):} A highly reliable alternative to MRV, particularly useful in emergency settings or for patients with contraindications to MRI.
    \item \textbf{Non-contrast brain CT scan:} Often the initial scan performed; it may be normal in up to 30\% of cases, but can show indirect signs such as the "cord sign" (hyperdense vein) or localised haemorrhage.
    \item \textbf{Coagulation and thrombophilia screen:} Detailed testing (including lupus anticoagulant, anticardiolipin antibodies, Factor V Leiden) to identify underlying hypercoagulable states.
    \item \textbf{Lumbar puncture:} Performed to measure cerebrospinal fluid opening pressure and rule out meningitis, but only after intracranial mass lesions or obstructive hydrocephalus are excluded by imaging.
\end{itemize}

\section*{Management}
The primary therapeutic goals are recanalisation of the occluded vein, prevention of thrombus propagation, and management of elevated intracranial pressure.
\begin{itemize}
    \item \textbf{Anticoagulation therapy:} Mainstay of treatment, initially with therapeutic low-molecular-weight heparin (LMWH) or unfractionated heparin, even in the presence of stable intracranial haemorrhage, followed by oral anticoagulants (warfarin or DOACs) for 3--12 months.
    \item \textbf{Management of raised intracranial pressure:} Use of osmotic therapy (mannitol or hypertonic saline), elevation of the head of the bed, and in refractory cases, decompressive craniectomy.
    \item \textbf{Antiseizure pharmacotherapy:} Initiation of anticonvulsants (e.g. levetiracetam) in patients who have experienced a seizure or have focal cortical lesions on imaging.
    \item \textbf{Septic thrombosis management:} Aggressive intravenous antibiotic therapy and surgical drainage of the primary infectious focus (e.g. mastoidectomy) if parameningeal infection is present.
    \item \textbf{Hydration and supportive care:} Maintaining adequate fluid balance to prevent dehydration, which can worsen venous stasis and thrombus progression.
\end{itemize}

\section*{Complications}
\begin{itemize}
    \item \textbf{Venous infarction and intracranial haemorrhage:} Ischaemia and subsequent rupture of congested vessels, leading to intracerebral bleeding.
    \item \textbf{Permanent neurological deficits:} Persistent motor weakness, sensory loss, cognitive impairment, or vision loss due to optic nerve atrophy from chronic papilloedema.
    \item \textbf{Status epilepticus:} Refractory, continuous seizure activity requiring intensive care unit admission and anaesthetic infusion.
    \item \textbf{Dural arteriovenous fistula:} Long-term vascular complication where abnormal direct connections develop between meningeal arteries and venous sinuses.
    \item \textbf{Herniation syndrome:} Brain displacement due to mass effect from severe cerebral oedema or large venous haematoma, leading to brainstem compression.
\end{itemize}

\section*{Prognosis and Outcomes}
The prognosis for meningeal vein thrombosis is generally favourable if diagnosed and treated early, with complete functional recovery in approximately 80\% of patients. This represents a significantly better prognosis than that of arterial ischaemic strokes. Mortality is low (under 5\% to 10\%) and is typically driven by severe underlying comorbidities (such as metastatic cancer or severe sepsis) or massive cerebral herniation. Recurrence rates are low, estimated at approximately 2\% to 5\%, and are minimised by appropriate duration of anticoagulation and management of hypercoagulable risk factors. Long-term outcomes are optimised by rehabilitation (occupational and physical therapy) and monitoring of vision and seizure control.

\section*{Prevention and Self-Care}
\begin{itemize}
    \item \textbf{Avoid combined hormonal contraceptives:} Using alternative, non-hormonal, or progesterone-only contraceptive options if there is a personal or family history of thrombosis.
    \item \textbf{Stay hydrated:} Drinking adequate water, particularly during illness, travel, or physical exertion, to prevent haemoconcentration.
    \item \textbf{Manage thrombotic risk during pregnancy:} Discussing prophylactic low-molecular-weight heparin with an obstetrician if known prothrombotic risks exist.
    \item \textbf{Promptly treat local infections:} Seeking medical care for ear, sinus, or dental infections to prevent intracranial spread.
    \item \textbf{Adhere to anticoagulation:} Taking prescribed blood thinners consistently and attending regular blood tests if monitoring is required.
\end{itemize}

\section*{Sources and Further Reading}
The following reputable organisations provide clinical guidelines and research on cerebral venous and meningeal thrombosis:
\begin{itemize}
    \item Stroke Foundation Australia. \textit{Clinical guidelines for stroke management and venous thrombosis.} https://strokefoundation.org.au/
    \item American Heart Association. \textit{Diagnosis and management of cerebral venous sinus thrombosis guidelines.} https://www.stroke.org/
    \item Mayo Clinic. \textit{Overview of cerebral venous thrombosis diagnostics and treatments.} https://www.mayoclinic.org/
    \item European Stroke Organisation. \textit{European consensus guidelines on cerebral venous thrombosis.} https://eso-stroke.org/
    \item Cleveland Clinic. \textit{Patient resources on brain blood clots and venous health.} https://my.clevelandclinic.org/
\end{itemize}
